\section{Empirical Analysis of Anisotropy in Language Model Representations}

The purpose of this section is to empirically characterize the geometric structure
induced by human language corpora in large language model representations.
Rather than proposing new algorithms, we focus on establishing a set of
data-driven facts about anisotropy that hold across models, layers, and training
settings.
In particular, we show that anisotropy consistently manifests as a separation
between a coherent, low-rank spectral spike and an incoherent, high-dimensional
long tail.
This empirical characterization provides the foundation for the averaging-based
spectral approximation developed in subsequent sections.

\subsection{Experimental Setup}

Our analysis is based on activations extracted from pretrained large language
models.
Unless otherwise stated, we analyze intermediate activation tensors from a
decoder-only transformer, using a batch of sequences drawn from a standard
language modeling corpus.
For a given layer, activations have shape $[B, L, H]$, where $B$ is the batch
size, $L$ is the sequence length, and $H$ is the hidden dimension.
We flatten the batch and sequence dimensions to obtain a token matrix
$X \in \mathbb{R}^{N \times H}$ with $N = B \cdot L$ tokens.

All analyses in this section are performed on raw activations without
preprocessing or normalization unless explicitly stated.
We emphasize that the observations reported below are qualitatively stable across
layers, batches, and model scales, and do not depend on architectural details or
training objectives.

\subsection{Spectral Signature: Spike and Long Tail}
\label{subsec:spike_long_tail}

We begin by examining the singular value spectrum of the token matrix
$b \in \mathbb{R}^{(BL)\times H}$ obtained by flattening the batch and sequence
dimensions of the activation tensor.
Let
\[
b = U \Sigma V^\top,
\qquad
\Sigma = \mathrm{diag}(\sigma_1,\sigma_2,\dots),
\]
denote its singular value decomposition.
Throughout this section, we define \emph{spectral energy} as the squared singular
values $\sigma_k^2$, so that the total energy equals the Frobenius norm
$\sum_k \sigma_k^2 = \|b\|_F^2$.

Figure~\ref{fig:spectrum_spike_tail_combined} summarizes the spectral signature
of anisotropy using the singular value spectrum and the cumulative spectral
energy distribution.
Across all layers and batches we analyzed, the spectrum exhibits a striking and
stable spike--tail structure: a small number of singular values dominate the
spectrum, followed by a broad and smoothly decaying long tail.

\begin{figure}[t]
  \centering
  \begin{subfigure}[t]{0.48\linewidth}
    \centering
    \includegraphics[width=\linewidth]{figures/spectrum_spike_tail.png}
    \caption{Singular value spectrum (log scale).}
    \label{fig:spectrum_spike_tail}
  \end{subfigure}
  \hfill
  \begin{subfigure}[t]{0.48\linewidth}
    \centering
    \includegraphics[width=\linewidth]{figures/spectrum_cum_energy.png}
    \caption{Cumulative spectral energy
    $\sum_{i\le k}\sigma_i^2 / \sum_i\sigma_i^2$.}
    \label{fig:spectrum_cum_energy}
  \end{subfigure}

  \caption{Spectral signature of anisotropy in language model activations.
  (\textbf{Left}) The singular value spectrum exhibits a sharp low-rank spike
  followed by a broad and smoothly decaying long tail.
  (\textbf{Right}) The cumulative spectral energy shows that a very small fraction
  of components captures a large share of the total energy, while hundreds to
  thousands of directions are required to explain the remaining variance,
  indicating a broad and information-rich tail.}
  \label{fig:spectrum_spike_tail_combined}
\end{figure}

Quantitatively, in the example shown, the top $1\%$ of singular components capture
more than $60\%$ of the total spectral energy.
Moreover, fewer than $20$ components suffice to explain $50\%$ of the energy,
while reaching $90\%$ and $95\%$ requires several hundred components
(e.g., $k\approx 390$ and $k\approx 799$, respectively).
These values are computed directly from the cumulative energy distribution shown
in Figure~\ref{fig:spectrum_cum_energy}, and are explicitly annotated in the plot.
At the same time, the singular value spectrum continues to decay smoothly over
thousands of directions, indicating that the long tail is broad and structured,
rather than negligible noise.

This spike--tail separation is robust across layers, model scales, and training
regimes, and has been independently observed in prior studies on low-bit
quantization and optimizer dynamics.
From a data-geometry perspective, these results establish anisotropy as a stable
and intrinsic property of language model representations.
They also motivate a key design principle for mitigating anisotropy: any
effective intervention must selectively suppress the dominant spike while
preserving the high-dimensional structure of the long tail, which encodes
semantically meaningful variation.

While this section establishes the existence and scale of the spectral spike, it
does not address whether this spike corresponds to shared structure or random
variance.
We address this question next.


\subsection{The Spike: Low-Rank and Coherent}
\label{subsec:spike_coherent}

We now examine the structure of the dominant spectral spike and show that it is
both low-rank and highly coherent across tokens.
This distinguishes the spike from random high-variance noise and provides the
key empirical foundation for averaging-based spike suppression.

\paragraph{Low-rank structure of the spike.}
Although the leading singular value is often emphasized, the spike is not
confined to a single direction.
Instead, its energy is distributed across a small number of leading singular
components, forming a low-dimensional subspace.
In the example analyzed here, the top singular component captures approximately
$32.5\%$ of the batch mean energy, while the top $5$ components already capture
more than $96\%$.
Projecting onto the top $10$ components captures over $99\%$ of the batch mean
energy, and projecting onto the top $1\%$ of singular directions ($k=41$)
recovers more than $99.9\%$.
This demonstrates that the spike is low-rank but not rank-one, and that the
shared structure is distributed across a small subspace rather than a single
singular direction.

\paragraph{Coherence with shared structure.}
Large singular values alone do not imply shared structure: a high-variance
direction could arise from uncorrelated fluctuations.
To directly test whether the spike corresponds to a coherent, token-shared
component, we compare the broadcasted column-mean matrix
$\mathrm{ColMean}_{\mu}(b)=\mathbf{1}\mu^\top$ with individual rank-$1$
reconstructions $\sigma_k u_k v_k^\top$ of the activation matrix $b$.
For each singular component $k$, we measure the absolute cosine similarity in
Frobenius inner product, which captures how strongly the shared mean aligns with
that component.

Figure~\ref{fig:colmean_alignment} shows the resulting alignment curves for three
different estimators of the column mean.
The horizontal axis is plotted on a logarithmic scale to emphasize the sharp
localization of alignment in the leading components.
Across all estimators, alignment is highly concentrated in the first few
singular directions and decays rapidly outside the spike subspace.
We note that the strongest alignment need not occur at $k=1$, reflecting the
fact that the spike occupies a low-rank subspace rather than a single direction.

\begin{figure}[t]
  \centering
  \begin{subfigure}[t]{0.32\linewidth}
    \centering
    \includegraphics[width=\linewidth]{figures/col_batch.png}
    \caption{Batch mean.}
    \label{fig:col_batch}
  \end{subfigure}
  \hfill
  \begin{subfigure}[t]{0.32\linewidth}
    \centering
    \includegraphics[width=\linewidth]{figures/col_sampled.png}
    \caption{Sampled tokens across sequences.}
    \label{fig:col_sampled}
  \end{subfigure}
  \hfill
  \begin{subfigure}[t]{0.32\linewidth}
    \centering
    \includegraphics[width=\linewidth]{figures/col_single.png}
    \caption{Single-sequence mean.}
    \label{fig:col_single}
  \end{subfigure}

  \caption{Alignment between the broadcasted column-mean matrix and individual
  rank-$1$ singular components under different mean estimators.
  The horizontal axis is shown on a logarithmic scale.
  Alignment is sharply localized to the leading singular components, indicating
  that the spectral spike corresponds to a coherent, shared structure.}
  \label{fig:colmean_alignment}
\end{figure}

\paragraph{Concentration of alignment mass.}
While the alignment curves reveal the shape and decay of coherence along the
spectrum, they do not directly quantify how much of the shared structure is
captured by the spike subspace.
To make this explicit, we compute the cumulative distribution function (CDF) of
alignment mass, defined as the fraction of total alignment captured by the first
$k$ singular components.

Figure~\ref{fig:colmean_alignment_cdf_individual} shows the CDFs for each mean
estimator separately.
In all cases, alignment mass is overwhelmingly concentrated in the spike
subspace.
For the batch mean estimator, more than $99.6\%$ of the alignment mass is
contained within the top $1\%$ of singular components.
The sampled-token estimator exhibits a similarly sharp concentration
($98.7\%$ within the top $1\%$), while the single-sequence mean is more diffuse
but still places over $98\%$ of the alignment mass in the same spike subspace.

\begin{figure}[t]
  \centering
  \begin{subfigure}[t]{0.32\linewidth}
    \centering
    \includegraphics[width=\linewidth]{figures/col_batch_cdf.png}
    \caption{Batch mean.}
    \label{fig:col_batch_cdf}
  \end{subfigure}
  \hfill
  \begin{subfigure}[t]{0.32\linewidth}
    \centering
    \includegraphics[width=\linewidth]{figures/col_sampled_cdf.png}
    \caption{Sampled tokens across sequences.}
    \label{fig:col_sampled_cdf}
  \end{subfigure}
  \hfill
  \begin{subfigure}[t]{0.32\linewidth}
    \centering
    \includegraphics[width=\linewidth]{figures/col_single_cdf.png}
    \caption{Single-sequence mean.}
    \label{fig:col_single_cdf}
  \end{subfigure}

  \caption{CDF of alignment mass for individual mean estimators.
  Each plot shows the fraction of total alignment captured by the top $k$
  singular components.
  In all cases, alignment mass is overwhelmingly concentrated in the spike
  subspace, with the batch mean exhibiting the strongest concentration and the
  single-sequence mean the most diffuse distribution.}
  \label{fig:colmean_alignment_cdf_individual}
\end{figure}

To facilitate direct comparison across estimators, Figure~\ref{fig:colmean_alignment_cdf}
overlays the three CDFs in a single plot.
This view highlights the consistent ordering:
batch mean $>$ sampled tokens $>$ single sequence, reflecting the role of
semantic diversity in reinforcing shared structure under averaging.

\begin{figure}[t]
  \centering
  \includegraphics[width=0.7\linewidth]{figures/colmean_alignment_cdf_three_methods.png}
  \caption{Comparison of alignment mass concentration across mean estimators.
  The batch mean captures over $99.6\%$ of the total alignment mass within the
  top $1\%$ of singular components, followed by the sampled-token estimator
  ($98.7\%$) and the single-sequence mean ($98.2\%$).
  This ordering reflects the role of semantic diversity in reinforcing shared
  structure under averaging.}
  \label{fig:colmean_alignment_cdf}
\end{figure}

\paragraph{Effect of averaging across sequences.}
Taken together, Figures~\ref{fig:colmean_alignment}--\ref{fig:colmean_alignment_cdf}
establish that the spectral spike is not only low-rank, but also coherent and
globally shared across tokens.
Averaging over the full batch yields the strongest and most concentrated
alignment, followed by averaging over randomly sampled tokens across sequences,
while single-sequence averaging is least effective.
This behavior directly supports the view that the spike corresponds to shared
linguistic structure, rather than uncorrelated variance, and that its coherence
is precisely what enables simple averaging to isolate it.

This property distinguishes the spike from the high-dimensional, incoherent long
tail examined in the next subsection.

\subsection{The Long Tail: High-Dimensional and Incoherent}
\label{subsec:tail_incoherent}

We next analyze the long tail of the spectrum and show that it exhibits the
complementary properties needed for safe spike suppression:
(i) semantic information is distributed over a high-dimensional subspace, yet
each token occupies only a portion of this subspace (compressibility); and
(ii) semantic components are largely incoherent across tokens drawn from
different sequences or sufficiently distant positions, so averaging does not
erase semantic variation.

\paragraph{Per-token compressibility in the tail eigenbasis.}
We define the spike subspace as the span of the top $k_s$ right singular vectors
(e.g., $k_s=1\%$) and project each token onto the tail complement.
We then express each projected token in the \emph{tail eigenbasis} (the remaining
right singular vectors) and measure how many tail coordinates are required to
capture a fixed fraction of its energy.
Figure~\ref{fig:tail_effective_support} shows the resulting distributions of
effective support sizes.

In our running example, the tail subspace has dimension $H-k_s=4055$.
A typical token requires only a few hundred tail coordinates to capture
$90\%$ of its tail energy (median $\approx 542$), and under $1000$ coordinates to
capture $95\%$ (median $\approx 896$).
Thus, although the semantic subspace itself is high-dimensional, individual
tokens are \emph{compressible} in this basis: each token activates only a subset
of the available tail directions.

\paragraph{Coefficient-space incoherence across tokens.}
Compressibility alone does not guarantee safe averaging: if different tokens
activate the same directions with aligned coefficients, semantic information
could still be suppressed.
What matters is therefore not overlap of supports, but \emph{incoherence of
coefficients} in the tail basis.

To test this directly, we compute cosine similarity between normalized tail
coefficient vectors for pairs of tokens, distinguishing three cases:
(i) tokens from the same sequence and nearby positions;
(ii) tokens from the same sequence but far apart in position; and
(iii) tokens drawn from different sequences.
Figure~\ref{fig:tail_coeff_cosine} shows the resulting distributions.

Across sequences, tail coefficient cosine similarity is tightly concentrated
around zero (mean $\approx -1.1\times 10^{-3}$, standard deviation $\approx 0.05$),
indicating near-orthogonality in coefficient space.
Tokens from the same sequence but far apart in position exhibit a similar
distribution (mean $\approx -4.4\times 10^{-3}$), suggesting that semantic
correlation decays with distance in long contexts.
In contrast, nearby tokens within the same sequence show a modest positive shift
(mean $\approx 0.076$), reflecting local semantic and topical coherence.

\begin{figure}[t]
  \centering
  \begin{subfigure}[t]{0.48\linewidth}
    \centering
    \includegraphics[width=\linewidth]{figures/tail_effective_support.png}
    \caption{Per-token effective support size in the tail eigenbasis.}
    \label{fig:tail_effective_support}
  \end{subfigure}
  \hfill
  \begin{subfigure}[t]{0.48\linewidth}
    \centering
    \includegraphics[width=\linewidth]{figures/tail_coeff_cosine.png}
    \caption{Cosine similarity of tail coefficients:
    within-sequence (near), within-sequence (far), and across-sequence.}
    \label{fig:tail_coeff_cosine}
  \end{subfigure}
  \caption{Empirical properties of the long tail.
  (\textbf{Left}) Individual tokens are compressible in the semantic tail: most
  tail energy is concentrated on a subset of directions.
  (\textbf{Right}) Tail coefficients are nearly orthogonal across sequences and
  for far-apart positions within a sequence, while nearby tokens show modest
  correlation, indicating local semantic structure but global incoherence.}
\end{figure}

Taken together, these results complement
Section~\ref{subsec:spike_coherent}.
The spike is low-rank and coherent, corresponding to shared linguistic
structure, whereas the long tail is high-dimensional and semantically diverse:
individual tokens occupy only a portion of the tail subspace, and tail
coefficients are weakly correlated across sequences.
This separation provides the empirical basis for the analysis in the next
section, where we formalize why averaging can suppress the coherent spike while
preserving semantic information in the long tail.

\subsection{Elementwise Origin of Outliers: Mean vs Spike vs Tail}
\label{subsec:outlier_origin}

The previous subsections characterized anisotropy in spectral space.
However, model training operates in \emph{element space}.
To understand why averaging is effective, we must identify which component
(mean, spike, or tail) actually produces elementwise outliers.

We decompose the activation matrix $X$ into three orthogonal components:
\[
X = M + X_{\text{spike}} + X_{\text{tail}},
\]
where:
\begin{itemize}
    \item $M = \mathbf{1}\mu^\top$ is the batch-mean component,
    \item $X_{\text{spike}}$ is the projection onto the top-$k$ singular vectors
          (e.g., top $1\%$),
    \item $X_{\text{tail}}$ is the remaining orthogonal complement.
\end{itemize}

We then analyze the largest $0.1\%$ of element magnitudes (quantization-relevant outliers)
and compute, for each outlier entry, the squared-magnitude contribution fraction from
mean, spike, and tail components.

\paragraph{Outlier-level contribution.}

Figure~\ref{fig:outlier_contribution} shows the distribution of squared-magnitude
shares within the outlier set.

\begin{itemize}
    \item For \textbf{activations}, the overwhelming majority of outliers are dominated
          by the \emph{mean component}. Over $99\%$ of extreme entries have their largest
          squared contribution coming from $M$, while spike contributes modestly and tail
          contributes negligibly.
    \item For \textbf{gradient outputs}, the mean contribution is negligible.
          Outliers are instead dominated by variance components (primarily tail).
\end{itemize}

This demonstrates a fundamental asymmetry:
activation outliers are primarily a \emph{mean-driven phenomenon},
whereas gradient-output outliers are variance-driven.

\begin{figure}[t]
    \centering
    \includegraphics[width=0.95\linewidth]{figures/outlier_contribution_hist.png}
    \caption{
    Squared-magnitude contribution shares for the top $0.1\%$ elementwise outliers.
    (Left) Activation matrix.
    (Right) Gradient-output matrix.
    Activation outliers are dominated by the mean component,
    while gradient-output outliers arise from variance (tail) components.
    }
    \label{fig:outlier_contribution}
\end{figure}

\paragraph{Global energy decomposition.}

While outliers are dominated by the mean, the overall Frobenius energy tells a
different story.
Figure~\ref{fig:energy_decomposition} shows the fraction of total energy
attributed to each component.

\begin{itemize}
    \item For activations, the mean and spike each contribute roughly
          $20\%$ of total energy, while the long tail contributes the majority.
    \item For gradient outputs, the mean contributes almost no energy;
          the spectrum is dominated by spike and tail variance.
\end{itemize}

Thus:
\begin{itemize}
    \item The \emph{mean dominates extreme elements} (outliers).
    \item The \emph{tail dominates total variance}.
\end{itemize}

This distinction explains why mean removal dramatically reduces outliers
while preserving most variance.

\begin{figure}[t]
    \centering
    \includegraphics[width=0.75\linewidth]{figures/energy_decomposition.png}
    \caption{
    Global Frobenius energy decomposition into mean, spike, and tail components.
    Activation anisotropy is driven by a mean component that creates
    elementwise outliers, whereas gradient-output anisotropy is primarily variance-based.
    }
    \label{fig:energy_decomposition}
\end{figure}

\paragraph{Key implication.}

The results clarify the mechanism underlying averaging-based conditioning:

\begin{itemize}
    \item Mean removal eliminates the dominant source of activation outliers.
    \item It does not significantly disturb the long-tail variance.
    \item Therefore, it improves quantization stability without destroying semantic information.
\end{itemize}

This elementwise perspective resolves the earlier confusion:
anisotropy in spectral space does not imply that spike variance dominates
element magnitudes.
Instead, coherent mean structure is responsible for most extreme values,

\subsection{Activation Bias and Optimizer Second-Moment Inflation}
\label{subsec:optimizer_bias}

We now quantify how activation bias affects adaptive optimizers.
Adam and AdamW maintain an elementwise second-moment estimate
\[
v \approx \mathbb{E}[g_t \odot g_t],
\]
where for a linear layer
\[
g_t = \delta_t x_t^\top.
\]
Using the identity
\[
v = \frac{1}{N}(\Delta^{\odot 2})^\top (X^{\odot 2}),
\]
we can directly measure how conditioning the activations $X$
affects the distribution of $v$.

\paragraph{Protocol.}
For each layer and training step, we compute:
\begin{itemize}
\item $v_{\text{raw}}$ from original activations,
\item $v_{\text{mean}}$ after mean removal,
\item $v_{\text{mean+whiten}}$ after mean removal and top-$1\%$ spike whitening.
\end{itemize}
We focus on extreme-tail statistics ($p99.9$ and $\max$), which dominate
effective learning-rate suppression in Adam updates.

\paragraph{Quantitative results.}

Table~\ref{tab:adam_v_detailed} reports raw values and reduction factors.

\begin{table}[t]
\centering
\caption{Extreme-tail statistics of Adam second moment $v$.}
\label{tab:adam_v_detailed}
\begin{tabular}{lcccccc}
\toprule
Layer & Step & Method & $p99.9(v)$ & $\max(v)$ &
$p99.9$ reduction & $\max$ reduction \\
\midrule
\multirow{3}{*}{Layer 3}
& 10k
& Raw
& $1.60\times10^{-10}$ & $8.38\times10^{-10}$ & -- & -- \\
& 
& Mean-only
& $3.44\times10^{-11}$ & $1.91\times10^{-10}$ &
$\mathbf{4.6\times}$ & $\mathbf{4.4\times}$ \\
&
& Mean+whiten
& $1.72\times10^{-11}$ & $8.09\times10^{-11}$ &
$9.3\times$ & $10.4\times$ \\
\midrule
\multirow{3}{*}{Layer 3}
& 180k
& Raw
& $1.00\times10^{-10}$ & $7.37\times10^{-9}$ & -- & -- \\
&
& Mean-only
& $2.50\times10^{-11}$ & $1.53\times10^{-10}$ &
$\mathbf{4.0\times}$ & $\mathbf{48\times}$ \\
&
& Mean+whiten
& $2.44\times10^{-11}$ & $1.54\times10^{-10}$ &
$4.1\times$ & $47.8\times$ \\
\midrule
\multirow{3}{*}{Layer 21}
& 10k
& Raw
& $1.67\times10^{-11}$ & $3.43\times10^{-11}$ & -- & -- \\
&
& Mean-only
& $1.59\times10^{-11}$ & $2.58\times10^{-11}$ &
$1.05\times$ & $1.33\times$ \\
&
& Mean+whiten
& $1.48\times10^{-11}$ & $2.33\times10^{-11}$ &
$1.13\times$ & $1.47\times$ \\
\midrule
\multirow{3}{*}{Layer 21}
& 180k
& Raw
& $9.23\times10^{-11}$ & $1.10\times10^{-9}$ & -- & -- \\
&
& Mean-only
& $2.21\times10^{-11}$ & $2.69\times10^{-10}$ &
$\mathbf{4.2\times}$ & $\mathbf{4.1\times}$ \\
&
& Mean+whiten
& $2.09\times10^{-11}$ & $1.95\times10^{-10}$ &
$4.4\times$ & $5.6\times$ \\
\bottomrule
\end{tabular}
\end{table}

\paragraph{Observations.}

\begin{enumerate}
\item Mean removal consistently reduces the extreme tail of $v$.
In early layers or later training stages, reductions range from
$4\times$ to nearly $50\times$ in $\max(v)$.

\item Whitening provides additional improvement in early training
(Layer 3, 10k), but contributes little beyond mean removal
in later stages (Layer 3, 180k).

\item The median of $v$ changes minimally across methods,
indicating that mean removal primarily affects a small subset
of extreme parameters, which nonetheless dominate adaptive
learning-rate suppression.
\end{enumerate}

\paragraph{Implication.}
These results demonstrate that activation bias is a dominant driver of
extreme second-moment inflation in Adam.
Removing the mean significantly stabilizes the optimizer by shrinking
the pathological tail of $v$, thereby improving effective learning-rate
allocation without altering the semantic covariance structure.